\documentclass[12pt]{article}
\usepackage[margin=0.72in]{geometry}
\usepackage{lmodern}
\usepackage{microtype}
\usepackage{multicol}
\usepackage{fancyhdr}
\usepackage{titlesec}
\usepackage{enumitem}
\usepackage{xcolor}
\usepackage[hidelinks]{hyperref}

\setlength{\columnsep}{0.30in}
\setlength{\parindent}{1.05em}
\setlength{\parskip}{0.24em}
\setlength{\headheight}{15pt}
\titleformat{\section}{\large\bfseries\scshape}{}{0pt}{}
\titleformat{\subsection}{\normalsize\bfseries}{}{0pt}{}
\pagestyle{fancy}
\fancyhf{}
\fancyhead[L]{Nietzsche on Truth and Metaphor}
\fancyhead[R]{Reading Handout}
\fancyfoot[C]{\thepage}
\renewcommand{\headrulewidth}{0.4pt}
\raggedbottom

\newcommand{\focusbox}[1]{%
\vspace{0.4em}\noindent\colorbox{gray!12}{\begin{minipage}{0.96\linewidth}#1\end{minipage}}\vspace{0.5em}}

\begin{document}
\begin{center}
{\LARGE\bfseries Do We Love Shadows?}\par\vspace{0.25em}
{\large\scshape Reading 6: Nietzsche, Truth, and Metaphor}\par\vspace{0.45em}
\rule{0.82\textwidth}{0.4pt}\par\vspace{0.7em}
\end{center}

\section*{Vocabulary Preview}
\begin{multicols}{2}
\begin{itemize}[leftmargin=*, itemsep=0.18em]
\item \textbf{Metaphor}: A comparison that carries meaning from one thing to another.
\item \textbf{Concept}: A general idea that groups many different things under one name.
\item \textbf{Convention}: A shared rule or habit a community agrees to follow.
\item \textbf{Anthropomorphic}: Shaped according to human needs, feelings, or perspective.
\item \textbf{Perception}: The way something appears to the senses or mind.
\item \textbf{Stimulus}: Something that affects the senses or nerves.
\item \textbf{Dissimulation}: Hiding, masking, pretending, or presenting oneself falsely.
\item \textbf{Tautology}: A statement that only repeats itself and adds nothing new.
\item \textbf{Abstraction}: A general idea formed by leaving out individual details.
\item \textbf{Subject}: The one who perceives, thinks, or speaks.
\item \textbf{Object}: The thing perceived, thought about, or named.
\item \textbf{Illusion}: An appearance mistaken for reality.
\end{itemize}
\end{multicols}

\section*{Introduction}
Friedrich Nietzsche lived from 1844 to 1900. In this early essay, written in 1873, he asks a disturbing question: when human beings say they possess truth, what are they actually possessing? His answer is meant to unsettle confidence. He argues that human beings make words, concepts, and systems that help them live, then forget that these truths began as human-made images.

This reading is the final complication in the unit. Plato asks us to turn from shadows to reality. Homer shows that illusions can be beautiful and knowledge-shaped. Bacon warns that the mind distorts truth. Augustine shows that love and habit can bind the will. Kant says enlightenment requires courage to use reason. Nietzsche now asks: \textbf{What if some things we call truth are metaphors we have forgotten are metaphors?}

\focusbox{\textbf{Source note:} Adapted and modernized from Friedrich Nietzsche, ``On Truth and Falsity in Their Ultramoral Sense'' (1873), in \textit{Early Greek Philosophy and Other Essays}, translated by Maximilian A. M\"ugge, edited by Oscar Levy, Project Gutenberg eBook 51548. This focused version preserves Nietzsche's opening fable, his account of truth as social convention, his language-and-concept argument, and his famous claim that truths are forgotten metaphors.}

\begin{multicols}{2}
\raggedcolumns
\section*{Reading}

\subsection*{1. A Fable About the Intellect}
In some remote corner of the universe, scattered among countless solar systems, there was once a star on which clever animals invented knowledge. That was the proudest and most deceitful minute in the history of the world, but it was only a minute. After nature had breathed for a short while, the star grew cold, and the clever animals had to die.

Someone could write a fable like this and still not show strongly enough how miserable, shadow-like, temporary, purposeless, and fanciful the human intellect appears within nature. There were eternities during which this intellect did not exist. When it disappears again, nothing will show that it ever existed.

The intellect has no mission beyond human life. It is purely human. Only its owner and maker treats it with such drama, as if the world turned around it. If we could speak with a gnat, we would learn that the gnat also flies through the air with great seriousness and feels itself to be the flying center of the world.

Nothing in nature is so small that, once touched by the power of knowledge, it does not swell up like a balloon. So the proudest human being, the philosopher, imagines the eyes of the universe are turned like telescopes toward his actions and thoughts.

\subsection*{2. Why Would Truth Arise Among Deceivers?}
The intellect was given to weak and short-lived creatures as a tool for staying alive. Its chief power is dissimulation: hiding, masking, and pretending. Creatures who cannot fight with horns or sharp teeth survive by deception.

In human beings this art reaches its highest point. Deception, flattery, lying, fraud, slander, display, disguise, acting before others and before oneself, and the constant fluttering around the flame of vanity are so common that one of the strangest questions is this: how could an honest impulse toward truth ever arise among human beings?

People are deeply immersed in illusions and dream-images. Their eyes slide over the surface of things and see forms. Their senses do not lead them to truth. They receive stimuli and play hide-and-seek on the backs of things.

At night, human beings let dreams lie to them their whole lives long, and their moral sense does not try to stop it. What does a person really know about himself? Nature hides most things from him, even about his own body: the winding intestines, the quick currents of blood, the delicate vibrations of nerves. Nature locks him inside proud, deceptive consciousness and throws away the key.

So where does the impulse to truth come from?

\subsection*{3. Truth Begins as a Social Agreement}
As long as individuals try to preserve themselves against other individuals, they usually use the intellect for deception. But because human beings need to live together, from necessity and boredom, they must make peace. They try to end, at least partly, the war of all against all.

This peace agreement brings something that looks like the first step toward truth. The community fixes what will count as ``truth.'' It invents names for things that are valid and binding for everyone. Language gives the first laws of truth.

Here the contrast between truth and lying begins. The liar uses the accepted names in order to make the unreal appear real. For example, he says, ``I am rich,'' when the correct name for his condition would be ``poor.'' He abuses the fixed conventions by switching terms.

If he does this selfishly and harmfully, society will no longer trust him. It will exclude him. In this way, people do not hate deception itself so much as they hate being injured by deception. Human beings desire truth in a limited sense: they want the pleasant, useful, life-preserving consequences of truth. They are often indifferent to pure knowledge that has no effect, and they may even be hostile to truths that would harm or destroy them.

\subsection*{4. Words Are Metaphors}
What are the conventions of language? Are they products of pure knowledge and love of truth? Do names and things perfectly match? Is language an adequate expression of reality?

Only by forgetfulness can a human being imagine that he possesses truth. If he will not be satisfied with tautology, with empty repetition, then he will always receive illusions instead of truth.

What is a word? It is the expression of a nerve-stimulus in sound. But to infer from that stimulus to a cause outside us is already a leap. How do we dare to say, ``the stone is hard,'' as if ``hard'' were something known in itself and not only a human sensation?

The thing-in-itself, the thing as it may be apart from our perception, is incomprehensible to the maker of language. He does not name the essence of things. He names only the relations of things to human beings. To express those relations, he uses the boldest metaphors.

First, a nerve-stimulus is transformed into a perception: first metaphor. Then the perception is copied into a sound: second metaphor. Each time there is a complete leap from one sphere into another.

When we speak of trees, colors, snow, and flowers, we believe we know something about the things themselves. Yet we possess only metaphors for things, metaphors that do not correspond directly to the original reality.

\subsection*{5. Concepts Flatten Difference}
Now think about how concepts are formed. Every word immediately becomes a concept, not because it perfectly remembers one original experience, but because it must fit countless similar cases. Yet similar really means never equal. Every concept arises by making unequal things equal.

No leaf is ever exactly the same as another leaf. Yet the concept ``leaf'' is formed by leaving out the individual differences and forgetting what makes each leaf distinct. Then we begin to imagine that besides all the particular leaves, there is something called ``the leaf,'' a true form according to which all leaves were made.

The same happens when we call a person ``honest.'' We ask, ``Why did he act honestly today?'' We answer, ``Because of his honesty.'' But this explains little. We have taken many different actions, left out their differences, named them honest actions, and then invented a hidden quality called honesty to explain them.

By disregarding the individual and real, we get the concept, just as we get the form. But nature knows no forms and concepts in this way. It knows only an inaccessible something that we cannot define.

\subsection*{6. What, Then, Is Truth?}
What, then, is truth? A mobile army of metaphors, metonymies, and anthropomorphisms: in short, a sum of human relations that have been poetically and rhetorically intensified, transferred, and decorated, and that after long use seem to a people fixed, official, and binding.

Truths are illusions we have forgotten are illusions. They are worn-out metaphors that have lost their sensory force. They are coins whose images have been rubbed away and that now count no longer as coins, but only as metal.

We still do not know where the impulse to truth comes from. So far we have heard only about the obligation society imposes in order to exist: to be truthful means to use the usual metaphors. Morally speaking, it is the obligation to lie according to a fixed convention, to lie with the herd in a style binding for all.

Of course, the human being forgets that this is what is happening. He lies in this way unconsciously, according to habits built up over centuries. By this very unconsciousness and forgetting, he arrives at a sense of truth.

As a rational being, he submits his actions to abstractions. He no longer lets himself be carried away only by sudden impressions. He builds a great structure of concepts, classifications, laws, ranks, and orders. This structure stands opposite the first world of perceptions and appears more fixed, more familiar, more human, and therefore more commanding.

\subsection*{7. Forgetting the Artist}
Human beings should be admired for this architecture of concepts. They build a complicated dome of ideas on a moving foundation, almost on running water. The material is fragile like cobwebs and yet firm enough not to be blown apart by every wind.

But this is not proof of a pure impulse toward truth. The seeker often finds only what he himself has hidden. If someone hides a thing behind a bush, looks for it, and finds it in the same place, there is not much to boast about.

Only by forgetting the original world of metaphors can a person live with calm, safety, and consistency. Only because a flowing mass of images has hardened into fixed concepts can he say, ``this sun, this window, this table is truth in itself.'' In short, only by forgetting that he himself is a subject, and an artistically creating subject, can he live securely.

If he could escape the prison walls of this faith even for one moment, his self-consciousness would be shattered. The question is not whether the bird's world, the insect's world, or the human world is the perfectly correct one. To decide that, we would need a standard of right perception that does not exist.

Between subject and object there is no perfect copying. At most there is an artistic relation, a suggestive transformation, a stammering translation into a completely different language.

\section*{Reading Focus}
\begin{enumerate}[leftmargin=*, itemsep=0.2em]
\item Underline Nietzsche's fable about the ``clever animals'' and knowledge.
\item Circle one sentence showing why human beings value truth for usefulness, not pure truth alone.
\item Put an M in the margin beside Nietzsche's claim that words are metaphors.
\item Put a C beside the section where concepts flatten different things into one general idea.
\item Box the sentence that best explains Nietzsche's answer to the unit question.
\end{enumerate}

\section*{Metaphor Map}
Choose one ordinary word, such as ``leaf,'' ``honest,'' ``success,'' ``smart,'' ``normal,'' or ``truth.'' Make a small map showing how that word gathers many different things under one name. Then write two sentences explaining what the word helps us see and what it makes us forget.

\section*{Final Comparison Claim}
Answer in one clear paragraph: \textbf{Does Nietzsche strengthen or challenge the whole unit?} State a claim, use one specific moment from Nietzsche, and compare him with Plato, Bacon, Kant, or another unit author.

\section*{Handout Summary}
Finish by writing a short summary of this handout. In three to five sentences, explain Nietzsche's account of language, concepts, and truth, and explain why this reading is a final complication for \textit{Do We Love Shadows?}

\end{multicols}
\end{document}
